\scanpage{6}
comme dans 2.4., pour tester \handwritten{l}a nullité de foncteurs comme dans2.1.
et dans 2.2. D'ailleurs\deleted{\unsure{r}}, 2.4. est évidemment u\handwritten{n} cas particulier de 2.5.
en prenant $\underline O_X$ comme faisceau ample sur le schéma affine $X$.

Soient encore $A$ un anneau noethérien, $\underline C$ la catégorie des $A$-modules\deleted{\unsure{,}}
de type fini, $\underline C'$ une catégorie abélienne quelconque, et
\[
T:\underline C\handwritten{\longrightarrow}\underline C'
\]
un foncteur additif. Pour tout $F\handwritten{\in}\mathrm{Ob}\,\underline C$, on peut alors considérer $T(F)$
comme un $A$-$\underline C'$-module, i.e. un objet de $\underline C'$ muni d'un\deleted{\unsure{e}} homomorphisme de
$A$ dans l'anneau de ses endomorphismes, De cette façon, le foncteur $F$ se
factorise à travers la catégorie des $\handwritten{A}$-$\underline C'$-modules, et on peut dire encore
que la donnée de $T$ équivaut à la donnée d'un \uline{foncteur} \inserted{$A$-}\uline{linéaire} de $\underline C$ dans
la catégorie \deleted{$\underline C'(A)$} $A(\underline C')$ des $A$-$\underline C'$-modules (plus précisément, \deleted{\unsure{en|ce}} on a une
équivalence naturelle entre ces deux catégories de foncteurs ...). C'est
le sorite déjà développé au par.7, mais qu'il faudrait faire de cette
façon plus générale, pour pouvoir par exemple renverser les flêches. Pour
l'instant, le fait qu'on s'est limité dans $\underline C$ aux modules de type fini
n'a pas joué, n'importe quelle sous-catégorie \inserted{pleine} de la catégorie des $A$-modules
aurait fait l'affaire (on se sert simplement du fait que les objets de $\underline C$
sont munis \deleted{``fonctoriellement} chacun d'une structure de $A$-$\underline C$-module, et que
les homomorphimses dans $\underline C$ sont des $A$-homomorphismes ...). Nous voulons
maintenant définir un homomorphisme fonctoriel
\[
T(A)\otimes_A M\handwritten{\longrightarrow}T(M)\tag*{\handwritten{$(*)$}}
\]
dont la \deleted{définition}\inserted{construction} utilisera le fait que \deleted{$A$ \unsure{$\in$} $\mathrm{Ob}\,\underline C$, et} $\underline C$ contient $A$ et que
les objets de $\underline C$ sont des modules de présentation finie sur $A$ (assuré par
$A$ noethérien, et les objets de $\underline C$ de type fini). Tout d'abord, si $\underline C'$ est
une catégorie abélienne, $A$ un anneau quelconque (pas néc comm.), \deleted{et} $\deleted{\unsure{L|M}}\handwritten{M}$ un
$A$-module à gauche, $P$ un $A$-$\underline C'$-module à droite, on définit
\[
P\otimes_A M\handwritten{\in\mathrm{Ob}}\,\underline C'\quad,
\]
sous réserve d'existence, comme l'\deleted{unique} objet de $\underline C'$ (déterminé à isomorphisme unique près) qui représente le foncteur \margin{\handwritten{déployer}}%
% slanted
$Q\handwritten{\leadsto}\operatorname{Hom}_A(M,\operatorname{Hom}(P,Q))$ de
$\underline C'$ dans $(\mathrm{Ens})$, -- où\handwritten{,} dans le double $\deleted{\unsure{H}}\operatorname{Hom}$, le $\operatorname{Hom}(P,Q)$ est considéré
comme un $A$-module à gauche, grâce aux opérations de $A$ sur $P$ à droite. Donc
par définition on doit avoir un homomorphisme fonctoriel en $Q$
\[
\operatorname{Hom}_{\underline C'}(P\otimes_A M,Q)\handwritten{\simeq}\operatorname{Hom}_A(M,\operatorname{Hom}(P,Q))
\]
% archive sheet number: 5

